\uuid{7ty4}
\exo7id{3630}
\titre{exo7 3630}
\auteur{quercia}
\organisation{exo7}
\datecreate{2010-03-10}
\isIndication{false}
\isCorrection{true}
\chapitre{Endomorphisme particulier}
\sousChapitre{Autre}
\cpge{mpsi-mp2i.dimension-finie.a, pcsi.dimension-finie.a, ptsi.dimension-finie.a, mpsi-mp2i.applications-lineaires.c, mpsi-mp2i.polynomes-fractions-rationnelles.a, pcsi.polynomes.a, ptsi.polynomes.a, bcpst1.espace-vectoriel-kn.b, bcpst2.applications-lineaires-matrices.b}
\module{Algèbre}
\niveau{L2}
\difficulte{}

\contenu{
\texte{
Soit $E = \R_3[X]$. On considère les formes linéaires :
$f_i : P  \mapsto  \int_{t=-1}^1 t^iP(t)\,d t$.
}
\begin{enumerate}
    \item \question{Montrer que $(f_0,f_1,f_2,f_3)$ est une base de $E^*$.}
    \item \question{Trouver la base duale.}
\reponse{
2. $\frac18(9-15X^2,\ 75X-105X^3,\ -15+45X^2,\ -105X+175X^3)$.
}
\end{enumerate}
}
